\renewcommand{\bibname}{Quellen und Bandzuordnung}
\begin{thebibliography}{99}
\addcontentsline{toc}{chapter}{Quellen und Bandzuordnung}
\footnotesize
\setlength{\itemsep}{6pt}
\setlength{\parskip}{0pt}
\interlinepenalty=10000

\bibitem{project}
Martin Kunze: \emph{Die Grundlagen der Mathematik}.
Band~45, insbesondere die Abschnitte über algebraische Halbverbände,
Verbände, Halbverbände mit Null und Eins sowie endliche Halbverbände
und Irreduzible; ergänzend \path{tex/b45/semilattice/statements.tex}.
Die Kapitel~1 und~7 dieses Leitfadens fassen wesentliche Kernaussagen
zusammen. Weitere Aussagen sind als bekannte Standardresultate ergänzt.
Das Format folgt den Leitfäden zu den Bänden~3, 4, 5, 28 und~37.

\bibitem{nation}
J. B. Nation: \emph{Notes on Lattice Theory}, Kapitel~2:
\emph{Semilattices, Lattices and Complete Lattices}.
Autorenfassung, insbesondere die algebraischen Definitionen,
Satz~2.2 zur Halbverbandsdarstellung sowie Lemma~2.8 und Satz~2.9 über
Irreduzible und Abschlußdarstellungen.
Hier werden nur die Halbverbands- und Verbandsresultate benutzt.
\url{https://math.hawaii.edu/~jb/math618/os2uh.pdf}.

\bibitem{smith}
Jonathan D. H. Smith: \emph{Key concepts for Boolean Algebra},
Lehrmaterial zur Veranstaltung M567, Iowa State University.
Einträge zu Halbverbänden, freien Halbverbänden und Idealen.
Die darüber hinausgehenden Strukturtheorien der Quelle sind
nicht Gegenstand dieses Leitfadens.
\url{https://jdhsmith.math.iastate.edu/class/M567Defn.htm}.

\bibitem{supclosed}
The mathlib Community: \emph{Mathlib.Order.SupClosed},
offizielle Dokumentation.
Abgleich von Unterverbänden und Erzeugung, insbesondere
\path{IsSublattice}, \path{latticeClosure},
\path{latticeClosure_empty}, \path{latticeClosure_singleton}
und \path{image_latticeClosure}.
\url{https://leanprover-community.github.io/mathlib4_docs/Mathlib/Order/SupClosed.html}.

\bibitem{nation-con}
J. B. Nation: \emph{Notes on Lattice Theory}, Kapitel~5:
\emph{Congruence Relations}.
Sätze~5.1 und~5.2 zu Kern und Quotienten sowie der Satz von
Funayama--Nakayama zur Distributivität des Kongruenzverbandes.
\url{https://math.hawaii.edu/~jb/math618/os5uh.pdf}.

\bibitem{freese-nation}
Ralph Freese und J. B. Nation:
\emph{Congruence lattices of semilattices}.
Pacific Journal of Mathematics 49 (1973), 51--58.
Vergleichsquelle zur Grenze der Übertragung von Verbandsresultaten
auf Halbverbände; mit Autorenkorrektur.
\url{https://msp.org/pjm/1973/49-1/pjm-v49-n1-p06-p.pdf};
\url{https://math.hawaii.edu/~ralph/Preprints/semicorr.pdf}.

\bibitem{nation-ideals}
J. B. Nation: \emph{Notes on Lattice Theory}, Kapitel~3:
\emph{Algebraic Lattices}, S.~28--29, Punkte~(4)--(6).
Begriffe der Halbverbands- und Verbandsideale, Erzeugungsformel
und ausdrückliche Unterscheidung der Signaturen mit und ohne Null.
Die weitere Theorie algebraischer vollständiger Verbände bleibt
außerhalb dieses Leitfadens.
\url{https://math.hawaii.edu/~jb/math618/os3uh.pdf}.

\bibitem{irreducible}
The mathlib Community: \emph{Mathlib.Order.Irreducible},
offizielle Dokumentation.
Abgleich der Begriffe \path{SupIrred}, \path{InfIrred} und der
Zerlegungssätze \path{exists_supIrred_decomposition},
\path{exists_infIrred_decomposition}.
Der Leitfaden benutzt die Aussagen für endliche Strukturen mit
Null beziehungsweise Eins.
\url{https://leanprover-community.github.io/mathlib4_docs/Mathlib/Order/Irreducible.html}.

\bibitem{nation-free}
J. B. Nation: \emph{Notes on Lattice Theory}, Kapitel~6:
\emph{Free Lattices}.
Satz~6.1 (Existenz), S.~64 (Größe), Satz~6.2 und das anschließende
Korollar (Whitman-Bedingung und Entscheidbarkeit).
Die Abzählbarkeit bei endlich vielen Erzeugern ist die unmittelbare
Folgerung aus der Darstellung durch endliche Terme.
\url{https://math.hawaii.edu/~jb/math618/os6uh.pdf}.

\bibitem{fjn}
Ralph Freese, Jaroslav Ježek und J. B. Nation:
\emph{Free Lattices}. Mathematical Surveys and Monographs~42,
American Mathematical Society, 1995.
Kapitel~I: Whitmans Lösung des Wortproblems.
\url{https://doi.org/10.1090/surv/042}.

\bibitem{dean}
Richard A. Dean:
\emph{Component subsets of the free lattice on n generators}.
Proceedings of the American Mathematical Society 7 (1956), 220--226.
Der Zusammenhang mit endlichen Gegenmodellen wird auch in
Nation, Kapitel~6, S.~71, unmittelbar vor Satz~6.11, beschrieben.
\url{https://math.hawaii.edu/~jb/math618/os6uh.pdf}.

\bibitem{freese-free}
Ralph Freese und J. B. Nation:
\emph{Free and Finitely Presented Lattices}, erweiterte Autorenfassung,
insbesondere Korollar~2-10.2 zur residualen Endlichkeit.
Für einen freien Verband beliebigen Rangs verwendet ein zu
trennendes Paar nur endlich viele Erzeuger.
\url{https://math.hawaii.edu/~ralph/Preprints/FreeLat-ExtendedVersion.pdf}.

\end{thebibliography}

\begingroup
\footnotesize
\noindent\textbf{Hinweise zum Quellengebrauch.}
Die Aussagen sind in der gemeinsamen Notation des Projekts
zusammengefaßt und auf ihre ausdrücklich genannten Voraussetzungen
abgestimmt. Die Quellen dienen dem Begriffs- und Satzabgleich;
die Formulierungen sind keine wörtlichen Übernahmen.
Literaturkonventionen zu leeren Trägern, leeren Idealen und
benannten Konstanten werden durch die Definitionen dieses Leitfadens
festgelegt. Die Online-Quellen wurden am 23.~Oktober 2026 eingesehen.
Eine historische Erstzuschreibung wird nur bei ausdrücklich
benannten Sätzen vorgenommen. Der Leitfaden enthält keine neuen
Forschungsergebnisse und beansprucht keine maschinelle Formalverifikation.
\endgroup
